% DN10-FRONTIER1
\section*{DN10-FRONTIER1: exact ten-clause frontier through seven variables}

\textbf{STATUS: COMPUTATIONAL / exact finite enumeration in the declared clause-presence model. P versus NP remains OPEN.}

Write \(M_n=8\binom n3\). Exact counting now gives
\[
U(4,10)=264{,}032,
\]
\[
U(5,10)=53{,}953{,}360,
\]
\[
U(6,10)=1{,}272{,}410{,}480,
\]
and
\[
\boxed{U(7,10)=14{,}141{,}034{,}880}.
\]
The corresponding ambient family sizes are \(\binom{M_n}{10}\), so the UNSAT fraction falls from approximately \(4.09\times10^{-3}\) at \(n=4\) to \(2.04\times10^{-8}\) at \(n=7\).

\paragraph{Full-support census.}
Let \(T_s\) denote the number of labeled UNSAT ten-clause formulas whose actual support is exactly a fixed set of \(s\) variables. Binomial inversion gives
\[
\boxed{T_4=264{,}032,\quad
T_5=52{,}633{,}200,\quad
T_6=952{,}650{,}800,\quad
T_7=6{,}357{,}940{,}960.}
\]
These values exactly reconstruct the ambient counts by
\[
U(n,10)=\sum_{s=4}^{n}\binom nsT_s
\]
for \(4\le n\le7\).

\paragraph{Genuinely minimal ten-clause cores.}
A private-assignment pruning rule was used: every clause of a minimally UNSAT formula must retain at least one falsified assignment not falsified by any other selected clause. Because private assignments can only disappear as clauses are added, loss of the last private assignment is a sound monotone pruning condition.

The resulting ambient minimal counts are
\[
29{,}792,\quad
18{,}229{,}720,\quad
342{,}221{,}760,\quad
2{,}862{,}493{,}200
\]
for \(n=4,5,6,7\). Their fixed-support core counts are
\[
\boxed{w_4=29{,}792,\quad
w_5=18{,}080{,}760,\quad
w_6=233{,}290{,}320,\quad
w_7=848{,}722{,}280.}
\]
The binomial reconstruction
\[
M_{10}(n)=\sum_{s=4}^{n}\binom ns w_s
\]
matches all four ambient minimal totals.

Thus the number of ten-clause UNSAT formulas that are not minimal is
\[
234{,}240,\quad
35{,}723{,}640,\quad
930{,}188{,}720,\quad
11{,}278{,}541{,}680
\]
for \(n=4,5,6,7\), respectively.

\paragraph{Independent checks.}
The \(D(6,10)\) total was obtained by fixed-clause symmetry, reproduced with reversed canonical candidate order, and reproduced again by pair-orbit double counting. The same pair-orbit implementation reproduces the established \(D(7,9)=60{,}047{,}120\) control. For \(D(7,10)\), a direct fixed-clause counter gives 505,036,960 UNSAT formulas containing a fixed clause and hence
\[
\frac{280}{10}\,505{,}036{,}960=14{,}141{,}034{,}880.
\]
An independent orbit-weighted pair-incidence calculation gives the identical total.

\paragraph{Lower-bound control.}
DN9-NEXT1 already proves that the certified all-pairs witness property persists across the complete nonconstant density range
\[
8\le m\le7\binom n3.
\]
Therefore the ten-clause cases satisfy
\[
B(n,10)\ge8\binom n3-1
\]
in the global clause-query DAG model. Numerically the certified lower bounds are
\[
31,\ 79,\ 159,\ 279
\]
for \(n=4,5,6,7\). For the declared connected fan-in-two AND/OR plus unary-NOT basis,
\[
G(n,10)\ge8\binom n3-2.
\]
The direct assignment-coverage construction still gives
\[
G(n,10)\le2^n\binom n3.
\]
These remain polynomial, representation-specific finite-model bounds and do not imply unrestricted SAT-time hardness.

\paragraph{Next structural target.}
For minimally UNSAT ten-clause formulas, Tarsi's deficiency bound limits actual support to at most nine variables. The next genuinely new fixed-support counts are therefore \(w_8\) and the deficiency-one endpoint \(w_9\). The latter is constrained by the classical structure of \(\mathrm{MU}(1)\); the former is the next computationally nontrivial case.

% END_OF_NOTE::DN10_FRONTIER1
